\documentclass[a4paper,11pt]{article}
\usepackage{exam-style-842-unit}
\examinehead{842 信息化概论 · 第 1 单元练习册}{练习一 参考答案与解析}
\begin{document}

\unitanswerscover{842 信息化概论 · 单元练习册}{第 1 单元参考答案}{练习一（概念与辨析）的答案与解析；每道判断题均给出改正后的正确表述。}

\answerpart{一、单项选择题（10 题，每题 3 分，共 30 分）}

\q B。\begin{analysis}数据是对客观事物的性质、状态和相互关系等进行记载的物理符号，是信息的载体，主要有数值、文字、声音及图形图像等形式（教材 1.1.1，第 1—2 页）。A 信息是数据加工后的内容；C 知识是使信息从定量到定性得以实现的、抽象的、具有逻辑的东西；D 智慧是从完全不同的知识中导出的一般性原理。\end{analysis}

\q A。\begin{analysis}香农认为“信息是能够用来消除不确定性的东西”（教材 1.1.1）。B 是维纳的观点：信息是我们适应外部世界、感知外部世界的过程中与外部世界进行交换的内容；C 是教材给出的一般定义，即信息是物质系统运动的本质特征；D 描述的是智慧。\end{analysis}

\q C。\begin{analysis}知识是让信息从定量到定性的过程得以实现的、抽象的、具有逻辑的东西，知识是为生产有意义的信息而生的（教材 1.1.1）。A 数据是记载物理符号；B 信息是加工后能减少不确定性的内容；D 智慧是从完全不同的知识中导出的一般性原理。\end{analysis}

\q B。\begin{analysis}二次信息是对一次信息进行加工处理后得到的信息，已经变成规则有序的信息，如文摘、索引、数据卡片，易于存储、检索、传递和使用（教材 1.1.1，第 2 页）。A 一次信息未经加工；C 三次信息是对一次、二次信息系统组织、压缩和分析的结果；D 动态型是按表现形式划分的类型（行情、战况等），与按加工程度分类不是同一维度。\end{analysis}

\q C。\begin{analysis}综述、专题报告、辞典、年鉴是系统地组织、压缩和分析一次信息和二次信息的结果，是人们深入研究的结晶，属于三次信息（教材 1.1.1）。A 一次信息是未经加工的原始信息；B 二次信息如文摘、索引；D 动态型信息强调时效性、生命周期很短。\end{analysis}

\q B。\begin{analysis}按表现形式划分，信息可分为文献型、档案型、统计型、动态型、图像型（教材 1.1.1）。A 按加工程度分为一次、二次、三次信息；C 按来源分为书本、报刊、电视、人、具体事物；D 时效长短不是教材使用的分类标准。\end{analysis}

\q D。\begin{analysis}信息的特点表现在 8 个方面：信息的普遍性、无限性和客观性，可共享性、可存储性、可传输性、可扩散性、可转换性、可度量性、可压缩性（教材 1.1.2，第 3 页）。可消耗性不在其中：信息具有可共享性，不会因使用而损耗。A、B、C 三项都是信息的基本特点。\end{analysis}

\q A。\begin{analysis}维系人类社会存在及发展的三大要素为物质、能源、信息（教材 1.1.2）。B、C、D 均不是教材表述；注意 D 中的信息、知识、智慧是认知层次，不是维系社会存在的三大要素。\end{analysis}

\q C。\begin{analysis}信息技术的四大范畴及其对应的人体器官是：传感技术是信息的采集技术，对应人的感觉器官；通信技术是信息的传递技术，对应人的神经系统；计算机技术是信息的处理和存储技术，对应人的思维器官；控制技术是信息的使用技术，对应人的效应器官（教材 1.1.4，第 4 页）。A 对应传感技术，B 对应通信技术，D 对应控制技术。\end{analysis}

\q B。\begin{analysis}一般把进行信息的收集、整理、存储、传输、处理及其应用服务的产业称为信息产业（教材 1.1.5）。A 只提到硬件制造，C 只提到通信线路设备，D 只提到软件，三者都只截取了信息产业的某一部分，范围过窄。\end{analysis}

\answerpart{二、多项选择题（4 题，每题 4 分，共 16 分）}

\q ABD。\begin{analysis}A 正确，教材 1.1.1 指出数据是信息的载体；B 正确，香农认为信息是能够用来消除不确定性的东西；D 正确，数据是计算机加工处理的对象，是未加工的对象。C 错误：数据是信息的载体，信息是数据经过加工以后能为某个目的使用的内容，是数据的内容或诠释，二者不能等同。\end{analysis}

\q ABC。\begin{analysis}信息的基本特点包括普遍性、无限性、客观性、可共享性、可存储性、可传输性、可扩散性、可转换性、可度量性、可压缩性（教材 1.1.2）。D 错误：可消耗性不是信息的特点，信息可共享，不会因使用而消耗。\end{analysis}

\q ABC。\begin{analysis}信息科学是研究信息现象及其运动规律和应用方法的科学，以信息论、控制论、系统论为理论基础，以电子计算机等为主要工具（教材 1.1.3，第 3—4 页）。D 错误：相对论是物理学的理论，不属于信息科学的理论基础。\end{analysis}

\q ABC。\begin{analysis}信息产业的特征包括：具有战略性的新兴主导产业，高渗透型、高催化型产业，知识、智力密集型产业，更新快、受科技影响大的变动型产业，高投入型产业，高增值型产业等（教材 1.1.5）。D 错误：与特征（6）“效益高的高增值型产业”正好相反。\end{analysis}

\answerpart{三、判断题（6 题，每题 2 分，共 12 分）}

\q √。\begin{analysis}教材 1.1.1 指出：数据是对客观事物的性质、状态和相互关系等进行记载的物理符号，是信息的载体；信息是数据经过加工以后能为某个目的使用的数据，是数据的内容或诠释。\end{analysis}

\q ×。改正：信息是数据经过加工以后能为某个目的使用的内容，是数据的内容或诠释；数据是信息的载体，二者不能等同。\begin{analysis}数据是计算机加工处理的对象，是未加工的对象；信息是加工后的结果。把信息与数据等同起来，混淆了载体与内容。\end{analysis}

\q ×。改正：系统论的奠基人是贝塔朗菲，控制论是维纳于 1948 年出版《控制论》一书创立的。\begin{analysis}维纳的研究对象是控制系统，控制论方法有功能模拟法和黑箱法；贝塔朗菲创立的系统论研究一切客观现实系统共同的特征、本质、原理和规律（教材 1.1.3，第 3—4 页）。\end{analysis}

\q √。\begin{analysis}狭义信息论是香农早期的研究成果，以编码理论为中心，主要研究信息系统模型、信息的度量、信息容量、编码理论及噪声理论；广义信息论又称信息科学，包括狭义信息论的内容，但研究范围比通信领域广泛得多（教材 1.1.3）。\end{analysis}

\q √。\begin{analysis}教材 1.1.3 的界定：科学是探知事物的本质、特征、内在规律以及与其他事物的联系，是关于自然、社会和思维的发展与变化规律的知识体系；技术是运用科学规律解决实现某一目的的手段和方法。工程则是将科学原理应用到工农业等生产部门中而形成的各门学科的总称。\end{analysis}

\q ×。改正：信息产业是需要大量智力和资金投入的高投入型产业，也是效益高的高增值型产业。\begin{analysis}信息产业特征（5）为高投入型产业，特征（6）为效益高的高增值型产业（教材 1.1.5），“低投入、低增值”与教材表述相反。\end{analysis}

\answerpart{四、配对题（6 对，每对 1 分，共 6 分）}

\q （1）—C\quad （2）—D\quad （3）—E\quad （4）—A\quad （5）—B\quad （6）—F。
\begin{analysis}香农是信息论创始人，1948 年发表《通信的数学原理》，1949 年发表《噪声下的通信》；维纳 1948 年出版《控制论》，是控制论创始人；贝塔朗菲是系统论奠基人；传感技术是信息的采集技术，对应人的感觉器官；计算机技术是信息的处理和存储技术，对应人的思维器官；综述、专题报告、辞典、年鉴属于三次信息。（对应关系见教材 1.1.3、1.1.4）\end{analysis}

\answerpart{五、简答题（3 题，每题 8 分，共 24 分）}

\q 参考答案要点：
（1）数据：对客观事物的性质、状态和相互关系等进行记载的物理符号，是信息的载体，主要有数值、文字、声音及图形图像等形式。
（2）信息：人类对自然世界的事物的变化和特征的反映，又是事物之间相互作用和联系的表征；信息是数据经过加工以后能为某个目的使用的内容。香农认为信息是能够用来消除不确定性的东西。
（3）知识：让信息从定量到定性的过程得以实现的、抽象的、具有逻辑的东西。
（4）智慧：从完全不同的知识中导出的一般性原理。
（5）相互关系：四者构成数据$\rightarrow$信息$\rightarrow$知识$\rightarrow$智慧的递进层次，抽象程度和价值逐级提高。数据是信息的载体，信息是数据的内容或诠释；信息经归纳、规律化上升为知识；知识经跨领域综合和情境判断上升为智慧；低层是高层的基础，高层是对低层的解释与运用。
（6）举例：田间测得“土壤含水率 18\%”是数据；结合作物需水规律判断“该地块需要灌溉”是信息；总结出“该品种在含水率低于 20\% 时产量明显下降”是知识；综合气象预报、市场价格、品种特性和灌溉成本，决定“今年是否调整种植结构与灌溉制度”则是智慧。

\q 参考答案要点：信息的特点是指信息区别于其他事物的本质属性，主要表现在 8 个方面（教材 1.1.2，第 3 页）：
（1）信息的普遍性、无限性和客观性；（2）信息的可共享性；（3）信息的可存储性；（4）信息的可传输性；（5）信息的可扩散性；（6）信息的可转换性；（7）信息的可度量性；（8）信息的可压缩性。
补充说明：可共享性指信息可被多个使用者同时使用而不损耗；可存储性和可传输性使信息能跨越时空积累与传播；可度量性使信息量可以用熵等指标度量；可压缩性指信息可通过压缩去除冗余。

\q 参考答案要点：
（1）信息技术主要研究信息的产生、获取、存储、传输、处理及其应用，是扩展人类信息器官功能的技术。
（2）四大范畴及对应器官：传感技术是信息的采集技术，对应人的感觉器官；通信技术是信息的传递技术，对应人的神经系统；计算机技术是信息的处理和存储技术，对应人的思维器官；控制技术是信息的使用技术，对应人的效应器官。
（3）信息科学是研究信息现象及其运动规律和应用方法的科学，以信息论、控制论、系统论为理论基础，以电子计算机等为主要工具，是一门新兴学科。
（4）信息技术与信息科学相互支撑：信息科学提供理论与方法，信息技术提供实现手段，二者共同推动信息产业发展。（教材 1.1.3、1.1.4）

\answerpart{六、综合题（1 题，共 12 分）}

\q 参考答案要点：
（1）归属判断（3 分）：田间观测记录是未经加工的原始信息，属于一次信息；病虫害简报是对一次信息加工处理后得到的规则有序信息，属于二次信息；年度专题报告是系统地组织、压缩和分析一次、二次信息的结果，属于三次信息。
（2）数据到信息的转化（3 分）：田间观测记录本身是零散的原始符号和数据，经整理、统计和诠释后成为能反映病虫害发生情况、减少不确定性的信息（简报）。
（3）信息到知识的转化（3 分）：把多次简报中反复出现的现象归纳为病虫害发生规律、防治指标和防治适期，形成规律化、逻辑化的知识。
（4）知识到智慧的转化（2 分）：综合多年专题报告、气象资料、市场价格与品种特性等不同来源的知识，导出一般性原理，作出是否防治、何时防治、如何轮作的决策，即智慧。
（5）结论（1 分）：一次信息是原料，二次信息提高了存储、检索、传递和使用的效率，三次信息是深入研究的结晶；三者逐级加工、不可偏废，最终为生产决策提供支持。

\end{document}
